\documentclass[11pt]{article}
\usepackage[T1]{fontenc}
\usepackage{amsmath,amssymb,amsthm}
\usepackage[margin=1in]{geometry}
\usepackage[hidelinks]{hyperref}
\newtheorem{theorem}{Theorem}
\newtheorem{lemma}[theorem]{Lemma}
\newtheorem{corollary}[theorem]{Corollary}
\newcommand{\m}{\mathfrak m}
\newcommand{\Cc}{\mathcal C}
\newcommand{\Dd}{\mathcal D}
\newcommand{\ord}{\operatorname{ord}}
\title{A same-centre conductor obstruction to geometric envelopes\\
on an irreducible pinch threefold}
\author{ROOT working paper 423: scoped progress on P12.3}
\date{9 October 2026}
\begin{document}
\maketitle

\begin{abstract}
Over every field of characteristic different from two, the irreducible
excellent local threefold
$k[u,x,y,s]_{(u,x,y,s)}/(x^2-y^2u)$ has a bounded strictly negative
continuous signed primary-model obstacle with no greatest geometric
continuous or decreasing-generated minorant. Its normalization is a
regular local ring with a single marked centre. The obstruction therefore
occurs within one normal centre, rather than between distinct components
or distinct marked centres. Explicit genuine primary ideals give positive
competitors on every semivaluation support. For an integer $L>2$, every
continuous minorant, including every canonical degree, has uniform error
at least $L-2$ from the raw supremum. The construction extends to every
dimension at least three. The two-dimensional pinch germ has continuous
envelopes at the qualified normal-surface scope. All geometric conclusions
use the entire qualified G331 comparison and descent theorem; the full
original problem and its deeper recursive source audit remain open.
\end{abstract}

\section{The geometric framework and the exact scope}
Let $\Cc_R$ be the geometric uniform completion of negative effective
finite real Cartier evaluations on the directed system of actual primary
blowups, with nonnegative degree on every contracted integral curve.
Let $\Dd_R$ consist of actual decreasing infima of $\Cc_R$ functions,
with the usual extended bottom convention. We use the entire qualified
G331 theorem: on the all-support normalized boundary $V_{\m}(R)$ these
cones agree with the corresponding ideal-profile completions; restriction
from finite normalization is the full quotient map, with its complete
conductor-infinite stratum; its positive descent and exact primary-gauge
transport hold. Every source, model, radial, support and historical-reading
qualification of that theorem is retained.

In particular, $\Cc_R$ functions are continuous and nonpositive,
$\Dd_R$ functions are upper semicontinuous, and an inequality valid for
every finite real negative primary-ideal profile passes first to
$\Cc_R$ by uniform limits and then to $\Dd_R$ by actual decreasing
infima. This substantive geometric comparison is the only incoming
correctness grant. The algebraic and all-support computations below are
proved directly. The standard finite-type excellence background is the
same as in the retained G331 and surface applications.

\section{A local irreducible normalization and its conductor}
Fix a field $k$ with $\operatorname{char}k\ne2$ and put
\[
 R=k[u,x,y,s]_{(u,x,y,s)}/(x^2-y^2u),\qquad
 \m=(u,x,y,s).
\]
The element $u$ is not a square in $k(u,y,s)$, so the defining quadratic
is irreducible and $R$ is a three-dimensional domain. It is excellent
as a local finite-type algebra over a field. In its fraction field put
$t=x/y$. Then $t^2=u$ and $x=yt$.

For an exact module computation, set
$A=k[u,y,s]_{(u,y,s)}$. The rings before the indicated localizations have
the respective $A$-module descriptions
\[
 R=A\oplus ytA,\qquad S=A\oplus tA,\qquad t^2=u.
\]
The finite integral $A$-algebra $S$ has just one maximal ideal over
$(u,y,s)$, namely $(t,y,s)$; it is therefore the local ring
$k[t,y,s]_{(t,y,s)}$. The ring $R$ also has one maximal ideal over that
ideal of $A$, so its displayed description is the stated localization.
The integral ring $S$ is regular and normal and has the same fraction
field as $R$. It is the finite normalization of $R$.

An element $r=a+ytb$ belongs to the conductor exactly when $rt\in R$:
since $S=R+Rt$, this is the entire conductor test. Now
$rt=at+yub$, and membership in $R=A\oplus ytA$ is equivalent to
$a\in yA$. Consequently
\[
 \mathfrak c=yS=(y,x)\subset R,\qquad
 R/\mathfrak c=k[u,s]_{(u,s)},\qquad
 S/\mathfrak c=k[t,s]_{(t,s)},\quad u=t^2.
\]
Both reduced conductor quotients are domains. The normalization has one
marked local centre. Its gauge extension is
\[
 \m S=(t^2,y,s).
\]

Define source semivaluations $\alpha,\beta$ by killing respectively
$(y,t-s)$ and $(y,t+s)$ and assigning value one to $s$ in the remaining
DVR. Both are normalized by $\m S$. They are distinct because $2$ is
invertible. Their restrictions coincide: each kills
$(y,x,u-s^2)$ and gives $s$ value one. Denote this target point by $o$.
It has nonminimal support. Both source points lie in the infinite
conductor stratum of the \emph{same} normal local centre.

\section{The signed obstacle and actual positive competitors}
Choose an integer $L>2$. The proper primary ideals
\[
 J_1=(x+ys)+\m^L,\qquad J_2=(y)+\m^L
\]
define the continuous signed primary-model obstacle
\begin{equation}\label{eq:F}
 F(v)=-v(\m)-v(J_1)+v(J_2)\qquad(v\in V_{\m}(R)).
\end{equation}
It is an actual signed real exceptional Cartier evaluation on a directed
common primary blowup. It is an obstacle, so its own positivity is not
required.

To verify its values on \emph{every} support, lift $v$ to a normalized
source semivaluation $w$ by G331 and set
$Y=w(y)$, $a=w(t+s)$. We allow $+\infty$. Since $w(\m S)=1$, we have
$Y\ge1$ and $a\ge0$. The product identity $x+ys=y(t+s)$ gives
\begin{equation}\label{eq:F-values}
 F(v)=-1-\{\min(Y+a,L)-\min(Y,L)\}.
\end{equation}
Only finite truncated quantities are subtracted; no infinity-minus-infinity
expression occurs. Thus $-L\le F\le-1$, including the entire
$Y=+\infty$ conductor stratum, where $F=-1$.

For $n\ge1$ the element
\[
 j_n=y(t+s)^n
\]
lies in $yS=\mathfrak c\subset R$. This is also explicit: expansion uses
$yt^{2i}=yu^i$ and $yt^{2i+1}=xu^i$. Define the genuine primary ideals
\[
 N_n=\m^{Ln}+(j_n),\qquad T_n=\m^nN_n,\qquad
 q_n(v)=-\frac{v(T_n)}n .
\]
Each contains a power of $\m$ and is proper. On every normalized lift,
\begin{equation}\label{eq:q-values}
 q_n(v)=-1-\min(Y/n+a,L).
\end{equation}
When $Y$ or $a$ is infinite this means the finite value $-1-L$.

For finite $Y,a\ge0$,
\[
 0\le\min(Y+a,L)-\min(Y,L)\le\min(a,L)
 \le\min(Y/n+a,L).
\]
The same inequalities hold when one value is infinite, by direct
truncation. Equations \eqref{eq:F-values} and \eqref{eq:q-values}
prove $q_n\le F$ on \emph{all} target supports.

The positivity of these competitors precedes any limit argument.
On the actual blowup $\operatorname{Bl}_{T_n}R$, the pulled-back
primary ideal is invertible and generated on each standard chart by a
nonzerodivisor, because $R$ is a domain and the chart embeds into its
fraction field. It defines an effective Cartier divisor $E_n$.
The line bundle $\mathcal O(-E_n)$ is relatively ample. Hence
$-E_n/n$ is a negative effective real Cartier presentation with
strictly positive degree on every contracted proper integral curve.
Therefore $q_n$ belongs to the actual geometric core of $\Cc_R$.

\section{A geometric upper bound at the identified point}
On $S$ use the regular coordinates $(s,z,y)$, where $z=t+s$.
Let $\nu$ be the generic monomial valuation of weights
\[
 \nu(s)=1,\qquad \nu(z)=L,\qquad \nu(y)=1 .
\]
It is centred at the same marked point and normalized by $\m S$:
$\nu(t)=\nu(z-s)=1$, so $\nu(t^2)=2$.
Write $\bar\nu$ for its restriction to $R$. Formula
\eqref{eq:F-values} gives
\begin{equation}\label{eq:nuF}
 F(\bar\nu)=-1-(L-1)=-L.
\end{equation}

For every proper primary ideal $I\subset R$,
\begin{equation}\label{eq:ideal-order}
 \bar\nu(I)\le o(I).
\end{equation}
Indeed $o(I)$ is finite, since $s^m\in I$ for some $m$.
Choose a generator attaining that value, and view it as a polynomial
in $(s,z,y)$ divided by a local unit. Its specialization $z=y=0$
is a nonzero polynomial in $s$ of order $e=o(I)$.
The nonzero $s^e$ monomial occurs in the numerator; the generic
monomial value is at most $e$, without cancellation between distinct
monomials. Both unit denominators have value zero. Unit ideals are
handled by value zero. This proves \eqref{eq:ideal-order} for all
ideals used in primary profiles, including all other semivaluation
supports involved in the geometric comparison.

Negating and summing \eqref{eq:ideal-order}, then applying G331's
\emph{entire} real geometric comparison and decreasing-infimum
extension, gives
\begin{equation}\label{eq:geometric-order}
 g(o)\le g(\bar\nu)\qquad(g\in\Cc_R\cup\Dd_R).
\end{equation}
Consequently every geometric minorant of $F$ satisfies
\begin{equation}\label{eq:upper}
 g(o)\le-L.
\end{equation}
This bound is proved for the full geometric cones, rather than for
a separately defined ideal-order class.

\section{All-element convergence and nonexistence}
For real $r\ge1$, let $w_r$ factor through
$S/(t-s)=k[y,s]_{(y,s)}$ and be the generic monomial valuation with
$w_r(y)=r$, $w_r(s)=1$. Let $v_r$ be its restriction to $R$.
It is normalized by $\m$. Since $\operatorname{char}k\ne2$,
$w_r(t+s)=w_r(2s)=1$. Thus
\[
 q_n(v_r)=-1-\min(r/n+1,L),\qquad
 \sup_{n\ge1}q_n(v_r)=-2.
\]
As $r\to+\infty$, $w_r$ converges to $\alpha$ on every element of
$S$: after substitution $t=s$, the monomial minimum of a polynomial
in $(y,s)$ tends to the order of its $y=0$ coefficient, or to infinity
if that coefficient vanishes. Local unit denominators have value
zero. Restriction therefore gives $v_r\to o$ on \emph{every}
element of $R$ in the actual pointwise boundary topology.

\begin{theorem}\label{th:failure}
The obstacle $F$ has no greatest $\Cc_R$ minorant and no greatest
$\Dd_R$ minorant. Both raw geometric suprema are finite and fail
upper semicontinuity at $o$.
\end{theorem}
\begin{proof}
A greatest minorant $g$ would dominate every $q_n$, so
$g(v_r)\ge-2$ for every $r$. Upper semicontinuity and $v_r\to o$
would force $g(o)\ge-2$. Equation \eqref{eq:upper} instead gives
$g(o)\le-L<-2$.

For either raw supremum $p$, the same competitors give $p(v_r)\ge-2$
and the universal bound gives $p(o)\le-L$. The constant $-1-L$
is a core competitor below $F$, so $p$ is finite. These inequalities
prove failure of upper semicontinuity directly.
\end{proof}

The two raw suprema agree. For completeness, if
$g=\inf_i g_i\in\Dd_R$ and $g\le F$, the directed ordinary open sets
$\{g_i<F+\varepsilon\}$ cover the compact boundary.
A single directed index covers it, and
$g_i-\varepsilon\in\Cc_R$ is below $F$. Taking pointwise suprema
and then $\varepsilon\downarrow0$ bounds $g$ by the raw
$\Cc_R$ supremum. The reverse inequality follows from
$\Cc_R\subset\Dd_R$.

\begin{corollary}\label{cor:gap}
Write $p$ for the common raw supremum. Every continuous geometric
minorant $g\le F$ satisfies
$\|p-g\|_\infty\ge L-2$.
In particular, for every $m\ge1$ the actual canonical ideal profile
\[
 A_m(F)=\{a\in R:v(a)+mF(v)\ge0\text{ for every }v\},\qquad
 h_m=-v(A_m(F))/m
\]
has $\|p-h_m\|_\infty\ge L-2$.
\end{corollary}
\begin{proof}
Along $v_r$, continuity gives
$g(v_r)\to g(o)\le-L$, whereas $p(v_r)\ge-2$.
Taking the supremum of the differences proves the inequality.
The canonical ideals contain $\m^{mL}$, are proper primary, and
their defining inequalities give $h_m\le F$ on every support.
Their blowups supply the same actual negative ample Cartier
presentations as above, so each $h_m$ is a continuous core minorant.
The proved every-degree raw-envelope convergence criterion from the
qualified canonical framework is consistent with this fixed positive
uniform gap: pointwise convergence cannot be uniform.
\end{proof}

The failure is stable. If $\|F'-F\|_\infty<\varepsilon$ with
$2\varepsilon<L-2$, the core minorants $q_n-\varepsilon$ of $F'$
force any greatest minorant to have values at least
$-2-\varepsilon$ along $v_r$. At $o$, \eqref{eq:geometric-order}
gives value at most $-L+\varepsilon$. Upper semicontinuity again
contradicts these bounds.

\section{Dimension, gauge and original residual}
For every $d\ge3$, adjoining $d-3$ variables to the displayed local
ring gives the same obstruction. Its normalization is the regular
local ring $k[t,y,s,z_1,\ldots,z_{d-3}]$ and its conductor is again
$yS$. Use the enlarged maximal ideal in $J_1,J_2,N_n,T_n$.
All-support formulas \eqref{eq:F-values}--\eqref{eq:q-values}
remain valid. Let $\alpha,\beta,w_r$ kill the extra variables,
and give each extra variable weight one in $\nu$.
The surviving $s^e$ monomial proves the identical universal
comparison. No higher-dimensional normal-source envelope theorem is
needed for any of these failures.

The two-dimensional pinch germ, obtained by omitting $s$, has reduced
conductor quotients $k[u]_{(u)}\subset k[t]_{(t)}$, $u=t^2$.
On their normalized boundary the only centred DVR valuation upstairs
has $v(t)=1/2$ when $v(u)=1$. The full normalization boundary map is
therefore bijective: the finite-conductor stratum is already singleton
by G331, and the conductor-infinite stratum has this unique lift.
Positive order/isometric descent and the qualified excellent-normal
surface theorem yield continuous envelopes for every continuous
signed obstacle. Adding one affine variable can thus destroy universal
CE while the normalization remains local and regular.

The characteristic-two construction is not covered: the two selected
lifts coincide. In that characteristic the reduced conductor quotient
map is purely inseparable, so G331 gives bijective boundary transport;
the higher-dimensional CE question becomes the corresponding regular
normal-source question at its actual characteristic hypotheses.

All functions in the construction are degree-one homogeneous on the
unnormalized cone. Under another primary gauge $\mathfrak b$, divide
by the same positive continuous $d(v)=v(\mathfrak b)$ and use the
exact G331 radial homeomorphism. The strict gap at $o$ becomes
$(L-2)/d(o)$, and the canonical defining ideals are unchanged for the
transported obstacle. The completion statement follows from the same
qualified all-support completion transport.

The construction supplies an irreducible \emph{nonnormal} failure
inside one normalization centre. General excellent normal
dimension-at-least-three envelope regularity, a universal criterion
for all conductor fibres, and the complete deeper recursive source
audit remain unresolved. No original P12.3 closure or new accepted
paper count is claimed. Full primary literal bodies remain LOCALONLY.
\end{document}
